\documentclass[../../../include/open-logic-section]{subfiles}

\begin{document}

\olfileid{sth}{story}{extensionality}
\olsection{विस्तारात्मकता}

सर्वांत आधी सांगायचे म्हणजे
संचांची ओळख त्यांच्या
!!{element}s वरून ठरते.
अधिक नेमकेपणाने:

\begin{axiom}[विस्तारात्मकता]
$A$ आणि~$B$ या संचांत
तेच !!{element}s
असतील, तर~$A$
आणि~$B$ हा
एकच संच आहे.
\[
  \lforall[A][\lforall[B][(\lforall[x][(x \in A \liff x \in B)] \lif
  \eq[A][B])]]
\]
\end{axiom}

\olref[sfr][][]{part}
या संपूर्ण भागात
आपण हे गृहीत धरले
होते. तरी ते
पुन्हा सांगण्यासारखे
आहे. विस्तारात्मकतेचे
स्वयंसिद्धक संचाची
ओळख त्याच्या
!!{element}s वरून
ठरते ही मूलभूत
कल्पना व्यक्त
करते. (म्हणून
संचांची तुलना
\emph{संकल्पनांशी}
करता येते:
अगदी त्याच
वस्तू अनेक
भिन्न संकल्पनांच्या
अंतर्गत येऊ
शकतात.)

हे तत्त्व का
स्वीकारावे? कारण
विस्तारात्मकता
नाकारणे म्हणजे
\emph{संच उपपत्ती}
म्हणता येईल
अशा गोष्टीचाच
त्याग करणे,
असे म्हणणे
विश्वसनीय वाटते.
संच उपपत्ती
म्हणजे विस्तारात्मक
संग्रहांचीच
उपपत्ती.

पण \olref[sth][][]{part}
मध्ये खरे आव्हान
\emph{कोणते संच
अस्तित्वात आहेत}
हे सांगणारी
तत्त्वे ठरवणे
आहे. या प्रश्नाचे
एकमेव खरोखर
‘‘उघड’’ वाटणारे
उत्तर सिद्धपणे
चुकीचे ठरते.

\end{document}
